Note que pelo teorema de Lagrange sabemos que para $n=1,2,3,5,7,\cdots,p$ com $p$ prima, temos que $G\cong \mathbb{Z}_{p}$.

Para $n=4$ note que pelo teorema de Lagrange se fixamos $x\in G$, temos dois casos $o(x)=4$, en ese caso $G\cong \mathbb{Z}_{4}$. Por outro lado, se $o(x)=2$, então tenemos el seguinte grupo
\begin{definition}[Grupo de Klein]
  Seja $K_4$ um grupo chamado el \textbf{grupo de Klein} caracterizado pela seguinte tabela de Cayley
  \[
  \begin{array}{c|cccc}
    \cdot & 1 & a & b & ab \\
    \hline
    1 & 1 & a & b & ab \\
    a & a & 1 & ab & b \\
    b & b & ab & 1 & a \\
    ab & ab & b & a & 1 \\
  \end{array}
  \]
\end{definition}
\begin{proposition}[Produto direito]
  Um grupo $G$ é o produto direito de dois subgrupos $H$ e $G$ se valen
  \begin{enumerate}
    \item $H\normal G$, $K\normal G$.
    \item $H\cap K=\{1\}$.
    \item $HK=G$.
  \end{enumerate}
\end{proposition}
\begin{notation}[Produto direito]
  Denotamos al grupo $G$ produto direito de $H$ e $K$ por $G=H\times K$.
\end{notation}
\begin{note}
  Note que $K_4\cong \mathbb{Z}_{2}\times \mathbb{Z}_{2}$, alem do mais
  \begin{center}
    \input{images/k4subgrupos}  
  \end{center}
\end{note}

Logo, para $n=6$ temos que si fixamos $x\in G$ com $x\neq 1$ que, se $o(x)=6$, então $G\cong \mathbb{Z}^{6}$.

Note que $o(x)=2$ no é possivel para  todo $x\in G$, dado que se for assim, então dados $a,b,c\in G$ existe $ab\in G$, logo tambem existe $ac$ e $cb$, pelo $|G|=6$, o que é uma contradição. Assim, existe $b\in G$ tal que $o(b)=3$.
Note que pelo teorema de Lagrange se $N=\langle b \rangle$, então $[G:N]=2$, por isso $N\normal G$, logo dado $a\in G$ tem-se que $aba^{-1}\in N$, logo
\begin{itemize}
  \item $aba^{-1}$ pode ser $1$?, não, dado que $o(b)=3$.
  \item $aba^{-1}$ pode ser $b$?, não, dado que $ab=ba$ e $o(ab)=6$.
  \item $aba^{-1}$ pode ser $b^{-1}$?, sim. 
\end{itemize}
Com isso $G\{1,a,b,b^{2},ab,ab^{-1}\}\cong S_3$.

Então temos
\begin{align*}
  |G|&=1\quad\text{ então }G\cong \{1\}.\\
  |G|&=2\quad\text{ então }G\cong \mathbb{Z}_{2}.\\
  |G|&=3\quad\text{ então }G\cong \mathbb{Z}_{3}.\\
  |G|&=4\quad,\text{ então } G\cong 
  \begin{cases}
    \mathbb{Z}_{4}, &\text{ se } o(x)=4 \text{,} \\
    K_4, &\text{ se } o(x)=2 .
  \end{cases}\, \text{ con }x\in G,x\neq 1.\\
  |G|&=5\quad\text{ então }G\cong \mathbb{Z}_{5}.\\
  |G|&=6\quad\text{ então }G\cong \begin{cases}
    \mathbb{Z}_6, &\text{ se } o(x)=6 \text{,} \\
    S_3, &\text{ se } o(x)\neq 6 .
  \end{cases},\text{ con }x\in G,x\neq 1.\\
  |G|&=7\quad\text{ então }G\cong \mathbb{Z}_{7}.\\
  |G|&=p\quad\text{ então }G\cong \mathbb{Z}_{p}.
\end{align*}
